% from ega3-6-fr.tex
\oldpage[III]{140}

\begin{corollary}[6.2.4]
Sous les hypothèses de (6.2.3), pour tout ouvert affine $V$ de $Y$,
l'homomorphisme canonique
\[
  \mathbf H^n(f^{-1}(V),\mathscr K^\bullet)
  \longrightarrow
  \Gamma(V,\mathscr H^n(f,\mathscr K^\bullet))
  \tag{6.2.4.1}\label{III.6.2.4.1.fr}
\]
est bijectif pour tout $n\in\mathbf Z$.
\end{corollary}

\begin{proof}
La démonstration est la même que celle de (1.4.11), en utilisant (6.2.2),
remplaçant $\mathscr F$ par $\mathscr K^\bullet$,
$\mathscr K^\bullet$ par
$f_*(C^\bullet(\mathfrak U,\mathscr K^\bullet))$,
$\mathscr H^\bullet(\mathscr K^\bullet)$ par
$\mathscr H^\bullet(f,\mathscr K^\bullet)$, et notant que ce dernier est un
$\mathscr O_Y$-Module quasi-cohérent par (6.2.3).
\end{proof}

\begin{proposition}[6.2.5]
Soient $Y$ un préschéma localement noethérien, $f:X\to Y$ un morphisme
propre, $\mathscr K^\bullet$ un complexe de $\mathscr O_X$-Modules tel que les
$\mathscr O_X$-Modules $\mathscr H^q(\mathscr K^\bullet)$ soient cohérents.
Alors les $\mathscr O_Y$-Modules $\mathscr H^n(f,\mathscr K^\bullet)$ sont
cohérents.
\end{proposition}

\begin{proof}
La question étant locale sur $Y$, on peut se borner au cas où $Y$ est
noethérien et affine, et il s'agit donc, en vertu de (6.2.4), de prouver que
les $\mathbf H^n(X,\mathscr K^\bullet)$ sont des
$\Gamma(Y,\mathscr O_Y)$-modules de type fini. On a alors
$\mathbf H^\bullet(X,\mathscr K^\bullet)=
\mathbf H^\bullet(\mathfrak U,\mathscr K^\bullet)$ (6.2.2), où l'on peut
supposer que $\mathfrak U$ est fini, puisque $X$ est quasi-compact. Les
cochaînes de chaque complexe $C^\bullet(\mathfrak U,\mathscr K^j)$ étant
alternées par définition, il y a un entier $r>0$ tel que
$C^i(\mathfrak U,\mathscr K^j)=0$ pour $i<0$ et $i>r$; on en conclut
(0, 11.3.3) que les deux suites spectrales du bicomplexe
$C^\bullet(\mathfrak U,\mathscr K^\bullet)$ sont birégulières. Comme les
intersections des ensembles de $\mathfrak U$ sont des ouverts affines
(I, 5.5.6), chaque foncteur
$\mathscr F\mapsto C^i(\mathfrak U,\mathscr F)$ est exact dans la catégorie
des $\mathscr O_X$-Modules quasi-cohérents; donc
$\mathrm H_I^q(C^i(\mathfrak U,\mathscr K^\bullet))
=C^i(\mathfrak U,\mathscr H^q(\mathscr K^\bullet))$, et les termes
$E_2$ de la seconde suite spectrale de
$C^\bullet(\mathfrak U,\mathscr K^\bullet)$ sont donnés (0, 11.3.2) par
\[
  ''E_2^{pq}
  =\mathrm H^p(C^\bullet(\mathfrak U,\mathscr H^q(\mathscr K^\bullet)))
  =\mathbf H^p(\mathfrak U,\mathscr H^q(\mathscr K^\bullet))
  =\mathrm H^p(X,\mathscr H^q(\mathscr K^\bullet))
\]
en vertu de (1.4.1); puisque $f$ est propre, ce sont des
$\Gamma(Y,\mathscr O_Y)$-modules de type fini (3.2.1). La suite spectrale
$''E(C^\bullet(\mathfrak U,\mathscr K^\bullet))$ étant birégulière, on en
déduit bien que les $\mathbf H^n(X,\mathscr K^\bullet)$ sont des
$\Gamma(Y,\mathscr O_Y)$-modules de type fini (0, 11.1.8).

A fortiori, si $\mathscr K^\bullet$ est un complexe de
$\mathscr O_X$-Modules cohérents, les $\mathscr O_Y$-Modules
$\mathscr H^n(f,\mathscr K^\bullet)$ sont cohérents sous les hypothèses de
(6.2.5) relatives à $Y$ et $f$ (0\textsubscript{I}, 5.3.4).
\end{proof}

\begin{env}[6.2.6]
L'hypercohomologie $\mathscr H^\bullet(f,\mathscr K^\bullet)$ est un foncteur
cohomologique dans la catégorie des complexes de $\mathscr O_X$-Modules
limités inférieurement (0, 12.4.4). C'est un foncteur cohomologique dans la
catégorie de tous les complexes de $\mathscr O_X$-Modules lorsque le
morphisme $f$ est quasi-compact et l'espace sous-jacent à $X$ localement
noethérien : en effet, il résulte alors de (G, II, 3.10.1) que $f_*$ permute
aux limites inductives (la question étant locale sur $Y$), et l'on peut
appliquer (0, 11.5.2).

Enfin, si $f$ est séparé, $\mathscr H^\bullet(f,\mathscr K^\bullet)$ est un
foncteur cohomologique dans la catégorie des complexes de
$\mathscr O_X$-Modules quasi-cohérents. C'est immédiat lorsque $Y$ est
affine, car alors $X$ est un schéma, donc, en vertu de l'isomorphisme
canonique (6.2.2), on est ramené à voir que
$\mathscr K^\bullet\mapsto\mathbf H^\bullet(\mathfrak U,\mathscr K^\bullet)$
est un foncteur cohomologique dans la catégorie des complexes de
$\mathscr O_X$-Modules quasi-cohérents, ce qui est immédiat puisque le
foncteur $\mathscr K^\bullet\mapsto
C^\bullet(\mathfrak U,\mathscr K^\bullet)$ est exact dans cette catégorie
(I, 1.3.7). Dans le cas général, pour tout ouvert affine $V$

